\section*{Введение}
\addcontentsline{toc}{section}{Введение}

\textbf{Актуальность темы.} Современный рынок розничной торговли электроникой характеризуется высоким уровнем конкуренции и постепенным 
переходом от классических интернет-магазинов к агрегаторам. Потребители все чаще предпочитают приобретать смартфоны, компьютерные комплектующие 
и периферию на платформах, объединяющих предложения от различных поставщиков, что позволяет сравнивать цены и выбирать оптимальные условия. Однако 
глобальные маркетплейсы зачастую перегружены товарами из других категорий, обладают сложной системой комиссий для продавцов и не 
предоставляют специализированных инструментов фильтрации, необходимых при выборе электронной техники.

В связи с этим возникает потребность в создании нишевых, профильных торговых платформ. Разработка специализированного маркетплейса электронной техники, 
предоставляющего малым и средним ритейлерам изолированные рабочие кабинеты, а администрации — прозрачные инструменты управления товарными потоками и 
аналитики, является актуальной задачей в сфере веб-разработки.

\textbf{Проблематика.} Существующие готовые системы управления контентом требуют значительных трудозатрат для адаптации под многопользовательскую 
архитектуру, где каждый продавец должен иметь независимый доступ к своим товарам и заказам без риска вмешательства в данные других участников. 
Использование современных веб-фреймворков позволяет решить эту проблему за счет проектирования гибкой ролевой модели, микросервисной архитектуры и 
оптимизированных реляционных баз данных.

\textbf{Объектом исследования} является организация процесса дистанционной торговли цифровой техникой в рамках многопользовательских платформ.

\textbf{Предмет исследования} — архитектурные подходы, алгоритмы агрегации данных и программные технологии, применяемые для разработки серверной и 
клиентской частей веб-маркетплейса.

\textbf{Цель курсового проекта} заключается в проектировании и программной реализации веб-платформы маркетплейса электронной техники, обеспечивающей 
взаимодействие продавцов и покупателей, а также сбор статистических данных для администратора.



